\section{Marking and representative laws}\label{mark:section}

This section prepares a consistent tuple for compression. A short marking
message first restricts most flags to domains of size $O(q^d)$. We then
choose a conditional law in which most positions have small entropy
deficit and small mutual information with a collection of representative
positions. In the reciprocal marginal classes, consistency turns this
information bound into a sparse-incidence estimate. All subspace dimensions
below are vector dimensions.

Fix $d$, $\eta$, and the scales from \eqref{base:parameters}. Let
$F=(F_1,\ldots,F_\ell)$ be a consistent tuple of deterministic length
$\ell=\Theta(k)$, extracted in either permitted orientation from the flag
stream. Assume that the stream law is dominated by a fixed constant times
the independent uniform law and that the simultaneous occupancy bound of
Lemma~\ref{base:rectangles} holds. An old context $\mathsf C$ is a discrete
random variable whose value specifies a domain for each slot, with
\begin{equation}\label{mark:stage-budget}
 H(\mathsf C)\le\Lambda,\qquad
 |\text{domain of }F_i\text{ given }\mathsf C|
       \le Cq^d e^\Delta,\qquad \Lambda,\Delta\ge0.
\end{equation}
The context and tuple may have arbitrary dependence on the stream.
Any public tables already fixed by preceding stages are regarded as
fixed here; further support and description tables will be drawn only
after the representative experiment below. Set
\begin{equation}\label{mark:parameters}
 D=\sigma^\beta(1+\Lambda/k+\Delta\sigma^{-\eta}),
 \qquad D\le\sigma^{1-\eta/2}.
\end{equation}
Thus $K=D\sigma^{3\beta}$ and $K_*=D\sigma^{6\beta}$, as in the common
parameter convention. Constants in this section may depend on $d,\eta$
and fixed retention and domination constants, but not on $q$ or on a
realized history.

\subsection{The marking message}

This marking procedure adapts the counting argument in
\cite[Claim~2.13]{Bradac2026}, with earlier expensive/cheap exposure
methods in \cite{AlonRodl2005}. The entropy bound below concerns the
actual conditional law of the selected flags.

In a forward scan, keep a vector subspace $V(y)$ of the second system for
each second point $y$. Initially it is zero. Whenever a position is
declared expensive, adjoin its second endpoint $b$ to $V(y)$ for every
$y$ incident with its first endpoint $a$. No other position changes the
state. Equivalently, $V(y)$ is the span of earlier expensive second
endpoints whose first endpoints annihilate $y$.

At a current flag $(a,b)$, consistency implies $a\perp V(b)$. Put
$r=\dim V(b)\le d$, and write $Z_j=\{y:\dim V(y)=j\}$. Choose
$l\in\{0,\ldots,r\}$ maximizing $|Z_l|$, with fixed tie-breaking.
Declare the position \emph{popular-cheap} if
\begin{equation}\label{mark:popular}
 |\{y\in Z_l:b\subseteq V(y)\}|\ge |Z_l|/(16q).
\end{equation}
If it is not popular-cheap, declare it \emph{poor-cheap} if
\begin{equation}\label{mark:poor}
 |\{y\in Z_l:a\perp y\}|\le |Z_l|/(8q).
\end{equation}
Declare it expensive otherwise, and update the state as described above.
Run the same scan on the reversed tuple with its two endpoint systems
interchanged. Consistency in this orientation follows from
Lemma~\ref{base:flag-graph}.
Let $E_+$ and $E_-$ be the expensive index sets in the two scans,
expressed in the original slot labels, and put $h=|E_+\cup E_-|$.

\begin{lemma}[Marking and entropy budget]\label{mark:marking}
Under \eqref{mark:stage-budget}, both scans together mark at most
$h\le Cq\sigma$ positions expensive. The old context, the two rank and
type masks, and the values of all positions expensive in either scan
form a message $\theta_0$ satisfying
\begin{equation}\label{mark:message-cost}
 H(\theta_0)\le\Lambda+O(k)
       +(d\sigma+\Delta+O(1))\E h.
\end{equation}
At a fixed $\theta_0$, every unspecified flag has a known joint domain
of size at most $e^{J_0}$, where $J_0=d\sigma+O(1)$ is deterministic.
For the set $U$ of unspecified slots, and every $\theta_0$-determined
subset $S\subseteq U$,
\begin{equation}\label{mark:subset-deficit}
 0\le \E\bigl[|S|J_0-H_{\theta_0}(F_S)\bigr]
       \le C(\Lambda+k+\Delta q\sigma).
\end{equation}
Here $H_{\theta_0}$ means entropy in the conditional law at the specified
value of $\theta_0$. Independently of the extraction entropy lower bound,
the same message gives
\begin{equation}\label{mark:entropy-upper}
 H(F)\le d\sigma\ell+\Lambda+O(k)+O(\Delta q\sigma).
\end{equation}
\end{lemma}

\begin{proof}
At a known state and rank $r$, a possible second endpoint $b$ has at
most $2q^{d-r}$ first partners annihilating $V(b)$. For $l\ge1$, counting
the incidences $b\subseteq V(y)$ over $y\in Z_l$ gives at most
$2|Z_l|q^{l-1}$ memberships. Hence \eqref{mark:popular} allows at most
$32q^l$ second endpoints, and its joint domain has size $O(q^d)$ since
$l\le r$. For $l=0$ there are no popular endpoints.

For a poor-cheap position, the variance estimate in
Lemma~\ref{base:incidence} bounds the number of first endpoints in
\eqref{mark:poor} by $Cq^{d+1}/|Z_l|$. Indeed their incidence count
deviates by at least $c|Z_l|/q$ from its mean for large $q$.
The second endpoint belongs to $Z_r$, and $|Z_r|\le|Z_l|$.
The product of these two marginal domain sizes is therefore at most
$Cq^{d+1}$. Incidence mixing bounds the number of flags between them by
\[
 Cq^{-1}q^{d+1}
       +q^{(d-1)/2}(Cq^{d+1})^{1/2}=O(q^d).
\]

At an expensive step, more than $|Z_l|/(8q)$ points of $Z_l$ are hit
by $a$, and fewer than $|Z_l|/(16q)$ already have $b\subseteq V(y)$.
Thus at least $|Z_l|/(16q)$ points increase their rank from $l$ to
$l+1$. The cumulative set
$U_l=\{y:\dim V(y)\le l\}$ has size at most $(d+1)|Z_l|$ at this
step, and loses a fraction at least $1/(16(d+1)q)$. Each $U_l$ is
nonincreasing throughout the scan. Starting with $O(q^d)$ points, a
nonempty integer-sized set can undergo only $O(q\sigma)$ such
multiplicative decreases. Charge a step to its chosen $l$ and sum over
$0\le l\le d$. This proves the expensive-position bound in each scan.

The constant-alphabet masks specify the ranks, types, and positions of
all transmitted flags. Their cost is $O(k)$; there is no separate cost
for the expensive positions. Conditional on the old context and masks,
the transmitted values range over a product of $h$ old domains.
Its logarithmic size is at most $h(d\sigma+\Delta+O(1))$, without any
independence assumption. This proves \eqref{mark:message-cost}.
The decoder reconstructs each scan by updating only at that scan's
expensive positions, ignoring values sent solely for the other scan.
Given the reconstructed state and the encoded rank $r$, the fixed
tie-breaking rule determines $l$; it need not be transmitted separately.
The decoder therefore recovers all cheap restrictions.
Choose $J_0=d\sigma+O(1)$ to bound their joint domains uniformly.

All values outside $U$ are known at $\theta_0$, so
\[
 \E H_{\theta_0}(F_U)=H(F\mid\theta_0)
                    \ge H(F)-H(\theta_0).
\]
Lemma~\ref{base:selection-entropy} and
\eqref{mark:message-cost} consequently give
\begin{align*}
 \E\bigl[(\ell-h)J_0-H_{\theta_0}(F_U)\bigr]
 &\le (\ell-\E h)J_0-H(F)+H(\theta_0)\\
 &\le \Lambda+O(k)+\Delta\E h-\eta\ell\log\sigma.
\end{align*}
The terms $d\sigma\E h$ cancel. Dropping the last, nonpositive term
gives \eqref{mark:subset-deficit} for $U$. For any specified history,
the chain rule and the coordinate caps imply
\[
 H_{\theta_0}(F_U)
 \le H_{\theta_0}(F_S)+|U\setminus S|J_0.
\]
Thus the deficit of $S$ is pointwise nonnegative and no larger than
that of $U$, proving the full assertion. Finally, the upper estimate
$H(F\mid\theta_0)\le(\ell-\E h)J_0$ combined with
\eqref{mark:message-cost} proves \eqref{mark:entropy-upper}.
\end{proof}

\subsection{Marginal classes and windows}

The same masks give more than a joint-domain bound. At a slot that is
poor-cheap in at least one scan, choose one such orientation by fixed
priority, and put
$u=\max\{\sigma,\min\{d\sigma,\log|Z_l|\}\}$ in that scan.
The sets of possible first and second endpoints, denoted by $A$ and
$B$, then satisfy
\begin{equation}\label{mark:reciprocal-caps}
 |A|\le A_{\max}=C e^{(d+1)\sigma-u},\qquad
 |B|\le B_{\max}=C e^u,\qquad \sigma\le u\le d\sigma.
\end{equation}
Before clamping, these are precisely the bounds
$Cq^{d+1}/|Z_l|$ and $|Z_l|$ proved above. Clamping at the lower end
uses the ambient first-space bound $O(q^d)$; at the upper end it uses
$|Z_l|\le Q_d=O(q^d)$. Hence it changes only the uniform constant.

If a slot is popular-cheap in both scans, write $r,r'$ for its forward
and reverse-swapped ranks. In the forward orientation,
\begin{equation}\label{mark:high-caps}
 \begin{gathered}
 |A|\le Cq^{r'},\qquad |B|\le Cq^r,\\
 \#\{a:(a,b)\text{ is allowed}\}\le2q^{d-r},\qquad
 \#\{b:(a,b)\text{ is allowed}\}\le2q^{d-r'}.
 \end{gathered}
\end{equation}
Here $1\le r,r'\le d$. If $r+r'\le d+1$, these bounds imply
\eqref{mark:reciprocal-caps} with $u=r\sigma$. Otherwise call the
slot a \emph{high slot} and classify it by its two ranks, keeping the
forward orientation. In a high class we use
$A_{\max}=Cq^{r'}$, $B_{\max}=Cq^r$.

Classify the reciprocal slots by orientation and by one of the bands
\[
 |u-r\sigma|\le K_*\quad(1\le r\le d),
 \qquad
 (r-1)\sigma+K_*<u<r\sigma-K_*\quad(2\le r\le d).
\]
Call these the integer and open bands, respectively. Since
$K_*=o(\sigma)$, they form a partition for large $q$, after assigning
boundary points to integer bands. There are only $O_d(1)$ classes.
Also $h=O(q\sigma)=o(k)$, so some class contains at least $ck$ slots
at every history, with a fixed sufficiently small $c>0$. Finite
pigeonholing supplies a fixed class and a $\theta_0$-determined event
of probability bounded below on which that class has at least $ck$
slots. Condition on this event, and retain a deterministic number
$\Theta(k)$ of its slots by a fixed $\theta_0$-based rule.

This conditioning leaves each law at a specified $\theta_0$ unchanged.
Because the subset deficit in \eqref{mark:subset-deficit} is
nonnegative, its expected value increases by at most a fixed factor.
Stream domination likewise changes only by a fixed factor. Reverse
the retained order and interchange endpoints when that is the class's
orientation; this preserves consistency and entropy. Relabel the
retained flags in this order. The domains and all their marginal
restrictions remain valid under any later conditioning.

We henceforth use the enlarged budget
\begin{equation}\label{mark:budget}
 \mathcal B=C(\Lambda+k+\Delta q\sigma),\qquad
 \mathcal B/k=O(D\sigma^{-\beta}).
\end{equation}
Discard an incomplete terminal window and partition the retained tuple
into $w$ windows of a common deterministic length $m$, divisible by four,
where
\begin{equation}\label{mark:windows}
 m\asymp
 \begin{cases}
 qD\sigma^{\eta/2},&\text{an open band},\\
 q\sigma^{1+\eta/2},&\text{an integer band},\\
 k,&\text{a high class, with }w=1.
 \end{cases}
 \qquad wm=\Theta(k).
\end{equation}
In the reciprocal cases $m=o(k)$, so this discarding loses $o(k)$ slots;
in the high case discard at most three slots. In each window take the
first and last quarters as representative blocks, leaving the middle
half for targets.

\subsection{A deterministic exposure round}

We now arrange that representative laws can be compared with most target
laws at a small information cost. It matters that the exposure round is
chosen from expectations, not after inspecting its flag values.

Conditioning to reduce dependence is a familiar correlation-rounding
principle; compare \cite[Lemma~4.5]{RaghavendraTan2011}. We require a
maximum over representatives from separate blocks, so we prove the
precise blockwise statement under the selected law.

Set $T_1=\lfloor qD\sigma^{3000\beta}\rfloor$. Before each of
$T_1$ rounds, independently in every representative block choose one
uniform unused index, independently also of all unexposed flags given
the preceding history; then expose the chosen flags. Write $\theta_t$
for the history before round $t$, starting with $\theta_0$, and $U_t$
for the unexposed positions. Each block then has
$n_t=m/4-t$ unused positions. The parameter inequalities give
\[
 \frac{T_1}{m}=O\!\left(
 \begin{cases}
 \sigma^{3000\beta-\eta/2},&\text{open},\\
 \sigma^{3000\beta-\eta},&\text{integer},\\
 \sigma^{3000\beta-3\eta/2},&\text{high},
 \end{cases}\right)=o(1).
\]
Thus $n_t=\Theta(m)$ throughout. Let $E_t$ be the fresh index set.
For $i\in U_t$, let $\mathcal J_i$ contain all members of $E_t$
except the representative in $i$'s own block, if $i$ is in a
representative block. At the fixed older history define
\[
 \delta_i=J_0-H_{\theta_t}(F_i),\qquad
 I_{ij}=I_{\theta_t}(F_i;F_j),\qquad
 M_i=\max_{j\in\mathcal J_i}I_{ij}.
\]
An empty maximum is zero. Fresh index draws do not change these
conditional flag laws.

\begin{lemma}[Representative exposure]\label{mark:representatives}
There is a deterministic round $t<T_1$, depending only on the current
joint law, such that, with $\theta=\theta_t$, $U=U_t$, and $E=E_t$,
\begin{equation}\label{mark:representative-bounds}
 \begin{aligned}
 \E\sum_{i\in U}\delta_i&\le\mathcal B,&
 \E\sum_{i\in U}M_i&\le C\mathcal B/T_1,\\
 \E\sum_{i\in E}\delta_i&\le C\mathcal B/m,&
 \E\sum_{i\in E}M_i&\le C\mathcal B/(mT_1).
 \end{aligned}
\end{equation}
All four expectations refer to the same pre-round history and fresh
indices. No values $F_E$ have yet been exposed.
\end{lemma}

\begin{proof}
At a history $h=\theta_t$, put
\[
 \mathcal D_t(h)=|U_t|J_0-H_h(F_{U_t}),\qquad
 \mathcal T_t(h)=\sum_{i\in U_t}H_h(F_i)-H_h(F_{U_t}).
\]
Both are nonnegative. The initial subset bound gives
$\E\mathcal D_0\le\mathcal B$, and
$\mathcal T_0\le\mathcal D_0$ since each marginal entropy is at most
$J_0$. For a fixed fresh index set $E_t$, the entropy chain rule gives
\[
 \mathcal D_t(h)
 -\E[\mathcal D_{t+1}\mid h,E_t]
       =|E_t|J_0-H_h(F_{E_t})\ge0.
\]
Consequently $\E\mathcal D_t\le\mathcal B$ at every round. Also
$\sum_i\delta_i\le\mathcal D_t$, proving the first bound at every
round, not just at the one eventually selected.

Expansion of the entropies gives the exact total-correlation identity
\begin{equation}\label{mark:tc-drop}
 \mathcal T_t(h)-\E[\mathcal T_{t+1}\mid h,E_t]
 =\operatorname{TC}_h(F_{E_t})+
        \sum_{i\in U_t\setminus E_t}I_h(F_i;F_{E_t}).
\end{equation}
Here $\operatorname{TC}_h(F_{E_t})$ is the sum of the selected marginal
entropies minus their joint entropy, at history $h$.
For each unselected $i$, data processing gives
$M_i\le I_h(F_i;F_{E_t})$. To include the selected positions, observe
that for an $i$ in a representative block, $M_i$ is a function of
the draws in the other blocks only. Its own block's draw is uniform
on $n_t$ indices and independent of those draws. Hence, conditional
on $h$ and averaging only over the fresh indices,
\begin{equation}\label{mark:selected-absorption}
 \E\left[\sum_{i\in E_t}M_i\,\middle|\,h\right]
 =\frac1{n_t}\E\left[
    \sum_{\substack{i\in U_t\\i\text{ in a representative block}}}
                      M_i\,\middle|\,h\right]
 \le\frac1{n_t}\E\left[\sum_{i\in U_t}M_i\,\middle|\,h\right].
\end{equation}
Combining this with \eqref{mark:tc-drop} yields
\[
 (1-n_t^{-1})\E\left[\sum_{i\in U_t}M_i\,\middle|\,h\right]
 \le \mathcal T_t(h)-\E[\mathcal T_{t+1}\mid h].
\]
Now average and sum over $t<T_1$. The right-hand side telescopes and
is at most $\E\mathcal T_0\le\mathcal B$, while $n_t\to\infty$
uniformly. Some deterministic $t$ therefore has
$\E\sum_iM_i\le C\mathcal B/T_1$.
Finally, uniform sampling gives
$\E\sum_{i\in E_t}\delta_i\le n_t^{-1}\E\sum_i\delta_i$;
\eqref{mark:selected-absorption} gives the identical factor for $M_i$.
Since $n_t=\Theta(m)$, both sampled bounds hold at that same round.
\end{proof}

\subsection{Endpoint mass bounds}

Fix the pre-round of Lemma~\ref{mark:representatives}. All definitions
in the rest of this section are made at its conditional law, before
revealing the fresh flags. Write
$p_i(a,b)$, $p_i^A(a)$, and $p_i^B(b)$ for the joint and endpoint
masses. A first endpoint $a$ is \emph{good} if
\[
 p_i^A(a)\ge e^{-K}/A_{\max,i},\qquad
 \Prob_\theta\bigl(p_i(a,B_i)>e^Kq^{-d}\mid A_i=a\bigr)\le .02.
\]
Define good second endpoints symmetrically, with $B_{\max,i}$.
For high classes impose also
$p_i^A(a)\le e^Kq^{-r'}$ and $p_i^B(b)\le e^Kq^{-r}$,
respectively. Endpoints outside the conditional support are not good.

\begin{lemma}[Endpoint mass]\label{mark:endpoint-mass}
At every fixed history and unexposed slot $i$, the probability that
either specified endpoint is not good is at most
\begin{equation}\label{mark:bad-endpoints}
 C\left(e^{-K}+\frac{\delta_i+1}{K}\right).
\end{equation}
This holds for both reciprocal and high classes, with their stated
marginal caps.
\end{lemma}

\begin{proof}
If a distribution $p$ is supported on at most $Ce^L$ atoms and
$Z=\log(p(X)e^L)$, then
$\Prob(-Z>t)\le\min\{1,Ce^{-t}\}$. Integrating this bound gives
$\E Z_-\le\log C+1$, after enlarging $C\ge1$. Thus
$\E Z_+=L-H(X)+O(1)$. Apply this with $L=d\sigma$ to the joint
law at slot $i$. The cheap-domain cap yields
\[
 \E_\theta[\log(p_i(F_i)q^d)]_+\le\delta_i+O(1).
\]
The mass of joint atoms exceeding $e^Kq^{-d}$ is at most
$C(\delta_i+1)/K$. By Markov's inequality, the marginal mass of
endpoints for which their conditional heavy-atom probability exceeds
$.02$ is at most fifty times this. The lower marginal cutoff excludes
mass at most $e^{-K}$ by the support-size cap.

For a high class, \eqref{mark:high-caps} also implies
\[
 H_\theta(A_i)\ge H_\theta(F_i)-(d-r')\sigma-O(1)
                  \ge r'\sigma-\delta_i-O(1).
\]
The same negative-part calculation, now using the support cap $Cq^{r'}$,
bounds the mass of $p_i^A(A_i)>e^Kq^{-r'}$ by
$C(\delta_i+1)/K$. The second marginal uses $r$ instead of $r'$.
\end{proof}

\subsection{Reciprocal cores and rectangle charges}

For reciprocal slots, we next convert small mutual information into
small incidence probability between good opposite endpoints. The
obstruction is a pair of flags in the same annihilator rectangle;
simultaneous stream occupancy will bound the total contribution of
such pairs.

Put $\rho=.01/(d+1)$. At a good first endpoint $a$ of slot $i$, define
a subspace in the second system by
\begin{equation}\label{mark:cores}
 V_i(a)=\bigcap_{\substack{x\text{ a first point}\\
       \Prob_\theta(B_i\not\perp x\mid A_i=a)\le\rho}}x^\perp.
\end{equation}
At a good second endpoint $y$ of slot $j$, define $U_j(y)$ in the first
system by the symmetric rule, using tests in the second system and the
law of $A_j$ given $B_j=y$. A basis chosen from the test covectors
uses at most $d+1$ tests, so each core captures at least $.99$ of its
conditional endpoint law. Incidence of every flag shows that $a$
itself is a zero-error test and likewise for $y$. Thus
\begin{equation}\label{mark:core-capture}
 \begin{gathered}
 \Prob_\theta(B_i\subseteq V_i(a)\mid A_i=a)\ge .99,
 \qquad a\perp V_i(a),\\
 \Prob_\theta(A_j\subseteq U_j(y)\mid B_j=y)\ge .99,
 \qquad y\perp U_j(y).
 \end{gathered}
\end{equation}
At a good $a$, joint nonheavy atoms have conditional mass at most
$q^{-d}A_{\max,i}e^{2K}$, and their intersection with the core has
mass at least $.97$. A vector subspace of dimension $v$ contains at
most $2q^{v-1}$ projective points. Since $K\to\infty$, this proves
\begin{equation}\label{mark:core-dimension}
 \begin{aligned}
 \dim V_i(a)&\ge
 \left\lceil d+1-\frac{\log A_{\max,i}+3K}{\sigma}\right\rceil,\\
 \dim U_j(y)&\ge
 \left\lceil d+1-\frac{\log B_{\max,j}+3K}{\sigma}\right\rceil.
 \end{aligned}
\end{equation}
The extra $K$ absorbs the constant from the factor $2/.97$.

\begin{lemma}[Reciprocal incidence and collisions]\label{mark:collision}
Fix a history $\theta$ in the selected pre-round. For unexposed slots
$i<j$ in a reciprocal class, write their flags as $(a,b)$ and $(c,y)$.
There are symmetric nonnegative numbers $c_{ij}=c_{ji}$, with
$c_{ii}=0$, such that
\begin{equation}\label{mark:incidence-charge}
 \Prob_{\mathrm{prod}}(a,y\text{ good},\ a\perp y)
                   \le C(I_{ij}+c_{ij}),
 \qquad \sum_{i<j}c_{ij}\le C\sigma k.
\end{equation}
Here the product law is the product of the two conditional flag
marginals at $\theta$. In an open band $c_{ij}=0$. In an integer band,
$c_{ij}$ is the conditional probability of a specified pair event
placing both flags in one rectangle. For high classes set $c_{ij}=0$;
no incidence estimate is asserted for them.
\end{lemma}

\begin{proof}
Under the actual conditional joint law, the event
$\{a\perp y,\ c\not\perp b\}$ has probability zero by consistency.
The relative entropy of this law to the product law is $I_{ij}$.
Lemma~\ref{base:event-transfer} therefore bounds the product
probability of this reverse-conflict event by $CI_{ij}$.
For given endpoints $a,y$, let
\[
 e(a,y)=\Prob_{\mathrm{prod}}(c\not\perp b\mid a,y).
\]
Except for good-incidence mass at most $CI_{ij}/\rho^2$, we have
$e(a,y)\le\rho^2$. For such endpoints, Markov's inequality says
that at least $1-\rho$ of the conditional $c$-mass satisfies
$\Prob_\theta(B_i\not\perp c\mid A_i=a)\le\rho$.
These $c$ are tests defining $V_i(a)$. Every projective point of
$V_i(a)$ consequently annihilates at least $1-\rho$ of the law of
$c$ given $y$, and is a test defining $U_j(y)$. Hence
\begin{equation}\label{mark:core-containment}
 U_j(y)\subseteq V_i(a)^\perp.
\end{equation}

To check the rounding even when $u_i\ne u_j$, take a common constant
$C\ge1$ in \eqref{mark:reciprocal-caps} and put $e_K=\log C+3K$.
Then \eqref{mark:core-dimension} reads
\[
 \dim V_i(a)\ge\left\lceil\frac{u_i-e_K}{\sigma}\right\rceil,
 \qquad
 \dim U_j(y)\ge\left\lceil d+1-\frac{u_j+e_K}{\sigma}\right\rceil.
\]
For large $q$, $e_K<K_*$ and $K_*+e_K<\sigma$.
In an open band indexed by $r$, we have
$(u_i-e_K)/\sigma>r-1$ and
$d+1-(u_j+e_K)/\sigma>d+1-r$.
Consequently the dimensions are at least $r$ and $d+2-r$.
These dimensions rule out
\eqref{mark:core-containment} in ambient vector dimension $d+1$.
All good-incidence mass is therefore exceptional, proving the claim
with $c_{ij}=0$.

In an integer band indexed by $r$, the two expressions inside the
ceilings are greater than $r-1$ and $d-r$, respectively, since
$u_i\ge r\sigma-K_*$ and $u_j\le r\sigma+K_*$.
The lower bounds are therefore $r$ and $d+1-r$. Containment forces
\begin{equation}\label{mark:complementary-cores}
 \dim V_i(a)=r,\qquad U_j(y)=V_i(a)^\perp.
\end{equation}
Let $\Gamma_{ij}$ be the endpoint event that $a,y$ are good and
incident and satisfy \eqref{mark:core-containment}. Define $G_{ij}$
by additionally requiring $b\subseteq V_i(a)$ and $c\subseteq U_j(y)$,
and put $c_{ij}=\Prob_\theta(G_{ij})$ in the actual joint law.
Given $a,y$ under the product law, $b$ and $c$ are independent with
their respective conditional endpoint laws. Therefore
\[
 \Prob_{\mathrm{prod}}(G_{ij})\ge .99^2
                 \Prob_{\mathrm{prod}}(\Gamma_{ij}).
\]
Applying Lemma~\ref{base:event-transfer} to $G_{ij}$ gives an upper
bound $C(I_{ij}+c_{ij})$ for its product probability. Add the
$CI_{ij}/\rho^2$ exceptional mass to obtain the first assertion
of \eqref{mark:incidence-charge}.

It remains to bound the sum of charges without treating the
history-dependent rectangles as fixed in advance. On $G_{ij}$,
\eqref{mark:core-capture} and \eqref{mark:complementary-cores} imply
\[
 b\subseteq V_i(a),\quad a\subseteq V_i(a)^\perp,
 \quad c\subseteq V_i(a)^\perp,\quad y\subseteq V_i(a).
\]
Both flags thus belong to $\mathcal R(V_i(a))$. Fix the history and
an actual realized tuple. For each fixed earlier slot $i$, all later
slots charged to it lie in the single rectangle determined by that
history and its actual first endpoint. The simultaneous occupancy
bound holds for every rectangle, including this adaptively chosen
one, on every stream in the support of the current law. No conditional
stream-domination estimate is needed. There are at most $C\sigma$
charged later slots for each $i$,
and hence at most $C\sigma k$ charged pairs, pointwise. Conditional
expectation proves the second assertion of
\eqref{mark:incidence-charge}.
\end{proof}

\subsection{Good positions and representatives}

The preceding information and collision bounds now give the uniform
input needed for the next compression step. Keep the fresh representative
indices, but still do not expose their values. The arrays $I_{ij}$ and
$c_{ij}$ are determined by the older history $\theta$ alone; only the
representative lists $\mathcal J_i$ depend on the fresh index draws.
Call an unexposed index
$i$ \emph{good} if
\begin{equation}\label{mark:good-index}
 \delta_i\le D\sigma^\beta,\qquad
 \max_{j\in\mathcal J_i}(I_{ij}+c_{ij})
                     \le \sigma^{-2000\beta}/q.
\end{equation}
This predicate depends only on the pre-round history and fresh indices.

\begin{proposition}[Good-index fractions]\label{mark:good-indices}
In every selected class, the expected proportion of bad unexposed
indices is $o(1)$, and so is the expected proportion of bad fresh
representatives. At every good index each endpoint is good with
probability $1-o(1)$, uniformly over the history and fresh indices.
All limits hold for fixed $d,\eta$ and the stage constants.
\end{proposition}

\begin{proof}
There are $\Theta(k)$ unexposed positions, because $T_1=o(m)$.
By \eqref{mark:budget} and Lemma~\ref{mark:representatives}, the
expected proportion violating the deficit bound is at most
\[
 C\frac{\mathcal B/k}{D\sigma^\beta}=O(\sigma^{-2\beta}).
\]
Using half the second threshold for mutual information, the corresponding
proportion is at most
\[
 C\frac{q\sigma^{2000\beta}\mathcal B}{kT_1}
                  =O(\sigma^{-1001\beta}).
\]
Only integer bands have nonzero collision charges. At a fixed history,
uniform sampling in blocks with $n_t=\Theta(m)$ gives
\[
 \E\left[\sum_{i\in U}\max_{j\in\mathcal J_i}c_{ij}
                       \,\middle|\,\theta\right]
 \le\frac{C}{m}\sum_{i,j\in U}c_{ij}
 \le\frac{C\sigma k}{m}.
\]
The expected proportion exceeding half the collision threshold is
therefore at most
\[
 Cq\sigma^{2000\beta}\frac{\sigma}{m}
                =O(\sigma^{2000\beta-\eta/2})=o(1).
\]
The maximum of $I_{ij}+c_{ij}$ is bounded by the sum of the two
separate maxima, proving the first assertion.

For fresh representatives, the expected deficit sum has an additional
factor $C/m$ by \eqref{mark:representative-bounds}. Each maximum over
the other blocks is independent of its own block's draw, so the same
selected-index identity \eqref{mark:selected-absorption} applies to
both mutual information and collision charges. Their expected sums
also gain a factor $C/m$. The number of fresh representatives is
$2w=\Theta(k/m)$, so their bad proportion obeys the same three
vanishing bounds. This is an average statement and does not require
that every representative be good simultaneously.

Finally, at a good index,
\[
 \frac{\delta_i+1}{K}
 \le \sigma^{-2\beta}+(D\sigma^{3\beta})^{-1}=o(1).
\]
Lemma~\ref{mark:endpoint-mass} proves the uniform endpoint assertion.
\end{proof}
